\documentclass[11pt]{article}
\usepackage[OT1]{fontenc}
\renewcommand{\rmdefault}{cmr}
\renewcommand{\sfdefault}{cmss}
\renewcommand{\ttdefault}{cmtt}
\usepackage{amsmath,amssymb,amsthm,mathtools}
\usepackage[margin=27mm]{geometry}
\usepackage{microtype,booktabs,longtable}
\usepackage[hidelinks]{hyperref}
\setlength{\emergencystretch}{2em}
\setlength{\parskip}{3pt plus 1pt}
\setcounter{tocdepth}{1}
\newtheorem{theorem}{Theorem}[section]
\newtheorem{lemma}[theorem]{Lemma}
\newtheorem{proposition}[theorem]{Proposition}
\newtheorem{corollary}[theorem]{Corollary}
\theoremstyle{definition}
\newtheorem{definition}[theorem]{Definition}
\newtheorem{example}[theorem]{Example}
\newtheorem{problem}[theorem]{Open problem}
\theoremstyle{remark}
\newtheorem{remark}[theorem]{Remark}
\newcommand{\N}{\mathbb N}
\newcommand{\NN}{\mathbb N_{>0}}
\newcommand{\Ftwo}{\mathbb F_2}
\newcommand{\lv}{\operatorname{lv}}
\newcommand{\norm}{\operatorname{N}}
\newcommand{\sz}{\operatorname{sz}}
\newcommand{\can}[1]{\mathopen{[\![}#1\mathclose{]\!]}}
\newcommand{\mle}{\preccurlyeq}
\newcommand{\mslt}{\prec}
\newcommand{\PhiF}{\Phi}
\newcommand{\Esrc}{\mathcal E_{\mathrm{src}}}
\newcommand{\headpart}{\operatorname{lim}}
\newcommand{\wt}{\operatorname{wt}}
\newcommand{\Ty}{\operatorname{Ty}}
\newcommand{\FV}{\operatorname{FV}}
\newcommand{\NF}{\operatorname{NF}}
\newcommand{\sgn}{\operatorname{sign}}
\newcommand{\id}{\operatorname{id}}
\newcommand{\BH}{\mathsf{BH}}
\newcommand{\ord}{\operatorname{ord}}
\newcommand{\bodyurl}[1]{\url{#1}}
\numberwithin{equation}{section}

\title{Controlled Ordinal Descents, Restricted Partitions\\and Cartier Observers}
\author{}
\date{Working manuscript, October 2026}
\begin{document}
\maketitle
\begin{abstract}
An ordinal descent can terminate at every input while its accumulated
step count retains difficult modular information. We give exact count
identities for a norm-controlled descent below $\varepsilon _0$. For a
polynomial block, congruent growing budgets give the same count modulo a
fixed modulus as a constant-budget restricted-partition problem. Applied
to a genuinely retained prefix in the Wilken--Weiermann assignment, the
identity produces a restricted-partition coefficient whose part bound
and selected index both vary along one prescribed fast-growing orbit.
We construct its complete finite coefficient observer. Differential
Cartier removes the higher poles, but on the actual linked inputs the
tuple of residues at one family of primitive roots has exact period
$2(\log _2 B-1)$. Thus a uniformly bounded cycle of the complete rational
state cannot solve this family. The result does not determine the
selected scalar coefficient, prove its nonperiodicity, or decide general
canonical comparison. Complete counts, counts to a boundary, full states,
and scalar observations are kept distinct throughout.
\end{abstract}
\tableofcontents

\section{Why a termination proof leaves a counting problem}

Wilken and Weiermann associate typed terms with ordinal vectors whose
zeroth coordinate is a natural number obtained from a norm-controlled
collapse \cite{ww}. An exact evaluation follows an ordinal descent. Its
ordinal argument decreases, but its available norm increases rapidly.
The number of replacements is part of the answer, rather than merely a
bound on the time needed to compute it.

This distinction matters for universal comparison. An algorithm that
eventually evaluates each supplied input need not decide whether
\emph{some} input has a specified residue. Nor does a large growth rate
determine a small residue. The methods below make that distinction
explicit by following one actual source boundary all the way to a
restricted-partition coefficient.

Our contribution has three levels. First, we prove an exact modular
transfer theorem for polynomial blocks under congruent growing budgets.
Second, we derive the correlated parameters of a block inside a
nonpolynomial source descent. Third, we analyze a complete finite
observer for its selected parity and prove an unbounded full-state
period on those very parameters. The third result limits a particular
compression strategy; it is not an impossibility theorem for all
methods of observing the final bit.

\section{Norm, control and the exact descent}

All ordinals in this paper lie below $\varepsilon _0$ and are given by
finite recursive Cantor normal forms. In displays of such forms,
addition and multiplication by finite coefficients have their ordinary
ordinal meaning. The norm is
\[
 \norm(0)=0,\qquad
 \norm\!\left(\sum_{i<r}\omega^{\alpha_i}c_i\right)
       =\sum_{i<r}c_i(1+\norm(\alpha_i)).
\]
Thus $\norm(n)=n$, $\norm(\omega^d)=d+1$ for finite $d$, and
$\norm(\omega^\omega)=3$. Every norm ball is finite and effectively
enumerable: its members are represented by finitely many rooted
expression trees of bounded size, after which their Cantor forms can
be compared effectively.

Fix exactly
\begin{equation}\label{eq:control}
 F_0(n)=2^n,\qquad F_{j+1}(n)=F_j^{\,n+1}(n),\qquad
 f(n)=\Phi(n)=F_5(n+100).
\end{equation}
Superscripts on functions denote iteration. In particular $f$ is
computable, strictly increasing, and
\begin{equation}\label{eq:growth}
 f(n)\ge 2^{n+100},\qquad f(n+1)\ge 2^{f(n)+1}.
\end{equation}
The finite-level iteration induction proving these growth estimates is
given in the foundations companion \cite{wwfoundations}; the modular
properties are proved below. The underlying collapse is
\[
 \psi(\alpha)=
 \max\bigl(\{0\}\cup
  \{\psi(\beta)+1:\beta<\alpha,\ \norm(\beta)\le f(\norm(\alpha))\}\bigr).
\]
The maximum is finite. It supplies the controlled implication
\begin{equation}\label{eq:controlled}
 \beta<\alpha,\quad \norm(\beta)\le f(\norm(\alpha))
       \quad\Longrightarrow\quad \psi(\beta)+1\le\psi(\alpha).
\end{equation}
Unrestricted monotonicity of $\psi$ is not asserted.

For $\alpha>0$, let
$M(\alpha,b)=\max\{\beta<\alpha:\norm(\beta)\le b\}$.
For a nonzero limit $\lambda$, its norm-saturating predecessor satisfies
\begin{equation}\label{eq:predecessor}
 \norm(M(\lambda,b))=b,\qquad
 \psi(\lambda+n)=\psi(M(\lambda,f(\norm(\lambda)+n)))+1.
\end{equation}
These identities, including the terminating Cantor-form algorithm for
$M$, are proved in the foundations companion
\cite{wwfoundations}. They are a specialization of the
norm-controlled predecessor framework of \cite{bcw}. In particular,
closed evaluation is an effective mathematical operation; no feasible
running time is implied.

\begin{definition}
For a limit $\lambda>0$ and total initial norm $t\ge\norm(\lambda)$,
start at $\alpha_0=\lambda+(t-\norm(\lambda))$. While $\alpha_i$ is
infinite, discard its finite tail and set
\[
 \alpha_{i+1}=M(\headpart(\alpha_i),f(\norm(\alpha_i))).
\]
Write $\tau_\lambda(t)$ for the number of replacements and
$Z_\lambda(t)$ for the terminal finite ordinal.
Set $J_\lambda(t)=\psi(\lambda+t-\norm(\lambda))$.
\end{definition}

\begin{proposition}[Complete-count identity]\label{prop:complete}
The descent terminates and
\[
 Z_\lambda(t)=f^{\,\tau_\lambda(t)}(t),\qquad
 J_\lambda(t)=Z_\lambda(t)+\tau_\lambda(t).
\]
\end{proposition}
\begin{proof}
Every replacement strictly decreases the ordinal, so there cannot be
infinitely many replacements. By norm saturation, its successive norms
are exactly $t,f(t),f^2(t),\ldots$. At the finite endpoint, the ordinal
itself equals its norm, proving the first identity. Each application of
\eqref{eq:predecessor} adds one. Finally $\psi(n)=n$ for finite $n$,
by induction in the defining maximum. Summing the contributions proves
the second identity.
\end{proof}

This count stops at the \emph{first finite ordinal}. It does not add a
further walk through the finite successors below $Z_\lambda(t)$.
Likewise, a count stopped at a nonfinite boundary is generally not
$\tau_\lambda$.

\section{The modular information supplied by the control}

We state the arithmetic input in a form that can be applied when a
modulus itself depends on a source parameter.

\begin{theorem}[Uniform residue envelopes]\label{thm:envelope}
There is an effective function $m\mapsto r_m\in\{0,\ldots,m-1\}$ such
that:
\begin{enumerate}
\item $f(x)\equiv r_m\pmod m$ for $x\ge m$;
\item for $s\ge1$, $j\ge2$ and $m\le2^{2^s}$,
      $F_j(s)\equiv r_m\pmod m$;
\item if $B=f(u)$, then
      $f(u+1)\equiv r_m\pmod m$ for every $m\le2^{2^B}$.
\end{enumerate}
\end{theorem}
\begin{proof}
We recall the construction to make the envelopes precise. Put $r_1=0$.
For $m>1$, write $m=2^a b$, $b$ odd, and set
$q=\varphi(b)$ if $b>1$, otherwise $q=1$. Then $q<m$.
Define $r_m$ by
$r_m\equiv0\pmod {2^a}$ and
$r_m\equiv2^{r_q}\pmod b$. This is a terminating recursion using
Euler's theorem and the Chinese remainder theorem.

Let $E_0(n)=n$, $E_{h+1}(n)=2^{E_h(n)}$.
Induction on $m$ proves that $E_h(n)\equiv r_m\pmod m$ whenever
$h,n\ge m$. The exponent is large enough for the $2$-part; its residue
modulo $q$ is supplied by the induction hypothesis, since $h-1\ge q$.
As $F_1(n)=E_{n+1}(n)$, every $F_j(n)$ with $j\ge1$ has residue
$r_m$ for $n\ge m$: in a higher-level iterate every argument is at
least $n$. This proves the first assertion.

For the stronger bound, the final $F_1$ call in
$F_2(s)=F_1^{s+1}(s)$ has argument at least
$F_1(s)\ge2^{2^s}$ when $s\ge1$. Its tower seed and height are both
at least $m$, proving the second assertion for $j=2$.
The induction to higher $j$ uses its last lower-level call, whose
argument is at least $s$.

Finally put $n=u+100$. In
$F_5(n+1)=F_4^{n+2}(n+1)$ the argument of the final $F_4$ is
$v=F_4^{n+1}(n+1)>F_4^{n+1}(n)=B$. The second assertion at $v$
therefore gives the third.
\end{proof}

The theorem does not say that an arbitrary computable sequence of
moduli may be stabilized at an arbitrary input. Each use must pay the
displayed bound. In particular, a lower bound on $f(x)$ by itself is
not the first assertion's lower bound on its \emph{argument}.

\section{Constant budgets and restricted partitions}

For $d\ge1$, define
\[
 p_d(n)=\#\{(a_1,\ldots,a_d)\in\N^d:
                         a_1+2a_2+\cdots+da_d=n\}.
\]
The independent choices of multiplicities give
\begin{equation}\label{eq:partition-series}
 \sum_{n\ge0}p_d(n)z^n=\prod_{i=1}^d(1-z^i)^{-1}.
\end{equation}
The relation to ordinals preserves every part, including the finite tail.

\begin{lemma}[A fixed-budget count]\label{lem:fixed}
Fix a nonnegative retained mass $c$ and take $S\ge c+d+1$.
Start the suffix at $\omega^d$ with any finite tail of total norm
(including $c$) at most $S$. At every replacement spend the constant
suffix budget $S-c$ and stop when the suffix is finite. The number of
replacements is $p_d(S-c)$. When $c=\norm(\pi)$ for a retained
CNF prefix $\pi$ of least exponent at least $d$, this is precisely
the corresponding segment starting at $\pi+\omega^d$ plus that tail.
\end{lemma}
\begin{proof}
The ordinals of suffix norm exactly $S-c$ below $\omega^d$ are
precisely
\[
 \omega^{d-1}a_d+\cdots+\omega a_2+a_1,\qquad
             \sum_{i=1}^d i a_i=S-c.
 \tag{*}
\]
Their number is $p_d(S-c)$. The first predecessor is the largest
such ordinal. In any later state of (*), the nonfinite coefficients
determine $a_1$ uniquely. Thus no second state of (*) has the same
nonfinite part. Removing the finite tail and taking the largest
predecessor of suffix norm $S-c$ selects exactly the next state of
(*) in decreasing ordinal order. All states are visited, ending at
the finite suffix $S-c$. There is one replacement into each state
and none after the last, proving the count. For $d=1$ there is one
state and one replacement.
\end{proof}

\begin{lemma}[A sufficient partition period]\label{lem:period}
For every $d,m\ge1$, put $Q_d=d!m^{d-1}$. Then
\[
             p_d(n+Q_d)\equiv p_d(n)\pmod m\quad(n\ge0).
\]
\end{lemma}
\begin{proof}
For $d=1$, $p_1(n)=1$. For $d>1$, splitting by the number
of parts of size $d$ gives
$p_d(n)=\sum_{j=0}^{\lfloor n/d\rfloor}p_{d-1}(n-dj)$.
Here $Q_d/d=mQ_{d-1}$. In the sum for $n+Q_d$, shifting the
index by $mQ_{d-1}$ recovers the old sum. The additional
$mQ_{d-1}$ summands form $m$ groups of $Q_{d-1}$ consecutive
indices in the periodic sequence
$j\mapsto p_{d-1}(n+Q_d-dj)\pmod m$.
Its period divides $Q_{d-1}$ by induction, so the extra sum is zero
modulo $m$.
\end{proof}

The period is deliberately a sufficient bound, not an optimal one.
For $d=2$ the count is $\lfloor n/2\rfloor+1$, and $2m$
is the period supplied by the lemma.

\section{A transfer theorem for congruent growing budgets}

For the numerical statement, a retained mass $c\ge0$ means that a
suffix of norm $n$ has total norm $c+n$, and a total budget $b$ leaves
suffix budget $b-c$. It need not be the norm of an ordinal prefix.
If it is realized by an actual retained CNF prefix, the numerical
process agrees with that source segment. A budget schedule may increase
with the current total norm; it must always afford that norm.

\begin{theorem}[Congruent-budget transfer]\label{thm:transfer}
Consider deletion of one suffix block $\omega^d$, with retained
numerical mass $c\ge0$ and initial total norm at least $c+d+1$,
under budgets at least the current total norm. Suppose
every budget used in this segment is congruent to one integer $S$
modulo $Q_d=d!m^{d-1}$, where $S\ge c$. Then its complete segment
count is congruent to $p_d(S-c)$ modulo $m$.
\end{theorem}
\begin{proof}
Induct on $d$. A degree-one block takes one replacement, as does
$p_1(S-c)$. For the induction step let $b$ be the first actual
total budget. The suffix of its greatest predecessor is the greedy
Cantor polynomial of norm $b-c$ and degree less than $d$.
Write its coefficients as $a_{d-1},\ldots,a_1$, with the final
finite coefficient included.

Thereafter each of its nonfinite monomial units is deleted by a
lower-degree block. For a degree-$i$ unit the retained prefix norm
has the form
\[
 c_i=c+\sum_{j>i}(j+1)a_j+(i+1)a,\qquad 0\le a<a_i.
 \tag{**}
\]
The order of those units is fixed by the Cantor order; growing
budgets affect their lower suffixes and durations, not the number
of these fixed higher monomial units. Since $Q_i\mid Q_d$, every
budget in a child segment has the residue of $b$ modulo $Q_i$.
Induction makes its count congruent to $p_i(b-c_i)$.
These arguments are nonnegative: each $c_i$ consists of $c$ plus
the norm of a prefix of the first suffix polynomial and is at most $b$.

Consequently the whole count, including its initial replacement,
has the same residue as the constant-$b$ process with exactly that
initial polynomial and those prefixes (**). By
Lemma~\ref{lem:fixed} this is $p_d(b-c)$.
Since $b\equiv S\pmod {Q_d}$, Lemma~\ref{lem:period} gives
$p_d(b-c)\equiv p_d(S-c)\pmod m$. If $b-c<S-c$,
apply the period lemma in the reverse direction by adding the
appropriate nonnegative multiple; no negative partition index is used.
\end{proof}

The proof covers the whole count, not just its leading contribution.
It does not replace the increasing operational budgets by a constant
control: the latter is only a device for computing one residue.

\begin{proposition}[An actual $\omega^\omega$ block]\label{prop:omega-block}
Let $\Omega=\omega^\omega$. At a real suffix state $\pi+\Omega+n$, with $\pi$ a retained CNF
prefix (zero or of least exponent at least $\omega$), and total norm $u$,
assume $c=\norm(\pi)$, $c+3\le u$, and
$B=f(u)\ge\max(16,m)$. Put $C=f(u+1)$ and $d=B-c-1$.
If $H_c(u)$ counts the complete deletion of this $\Omega$ copy,
then
\begin{equation}\label{eq:omega-block}
 H_c(u)\equiv1+p_{B-c-1}(C-c)\pmod m.
\end{equation}
Its output norm is still exactly $f^{H_c(u)}(u)$.
\end{proposition}
\begin{proof}
The first predecessor is $\pi+\omega^d$, of total norm $B$.
Indeed the greatest exponent of norm at most $B-c-1$ is the
finite exponent $d$ in this particular bound below $\omega$;
this is the predecessor of $\omega^\omega$ under the suffix budget.
All subsequent norms are at least $B$, and their budgets have
residue $r_{Q_d}$ by Theorem~\ref{thm:envelope}, because
\[
 Q_d=d!m^{d-1}\le B^{2B}\le2^{2^B}.
\]
The last inequality holds for $B\ge16$; it follows from
$2B\log_2 B\le2^B$, first at $16$ and then by the increasing
ratio of the right side to the left side.
The adjacent value $C=f(u+1)$ has the same residue modulo $Q_d$
by the third part of Theorem~\ref{thm:envelope}.
Apply Theorem~\ref{thm:transfer} to the remaining $\omega^d$ block
with $S=C$, then include its preceding $\Omega$ replacement.
Norm saturation proves the final assertion independently of the
modular calculation.
\end{proof}

\section{The first moving-prefix boundary}

Write $\Theta=\omega^{\omega+1}$. Suppose a genuine descent first
reaches a state whose limit part is $\Theta$, at total norm $v$.
The forthcoming calculation is conditional only on this actual
entry; it does not posit independent entry values for a larger source.
Put $u=f(v)$. The residue recursion gives $u\equiv1\pmod3$.

\begin{lemma}[The retained prefix at first entry]\label{lem:entry}
The predecessor of this $\Theta$ state is
\[
 \Omega q+1,\qquad q=(u-1)/3.
\]
The first $\Omega$ copy deleted in this polynomial has retained
prefix norm $e=u-4$ and total input norm $u$.
\end{lemma}
\begin{proof}
A copy of $\Omega$ has norm three. The maximal exponent below
$\omega+1$ that can occur is $\omega$; the remaining budget
is therefore spent on as many copies of $\Omega$ as possible and
then a finite remainder. Since $u\equiv1\pmod3$, that remainder
is one. Before deleting the last displayed $\Omega$ in the
ordinary sum $\Omega(q-1)+(\Omega+1)$, the retained prefix is
$\Omega(q-1)$ and has norm $3(q-1)=u-4$.
\end{proof}

With $B=f(u)$, $C=f(u+1)$, $e=u-4$ and $d=B-u+3$, the
first block parity is consequently
\begin{equation}\label{eq:delta}
 \delta_4(v):=H_{u-4}(u)\bmod2
       =1+p_d(C-e)\pmod2.
\end{equation}
The subscript records the norm four of its suffix $\Omega+1$.
The initial $\Theta$ replacement lies outside $H_{u-4}(u)$;
the displayed $1$ instead counts the $\Omega$ replacement.

This quantity is inside a larger correlated computation. Before its
remaining limit part has become exactly $\Omega$, the $\Theta$ segment
deletes $q-1$ copies with prefix masses $3(q-1),3(q-2),\ldots,3$. Their inputs
are successive points on the \emph{same} $f$ orbit. Observing the
sum of all these counts does not automatically define the graph of
the first summand. Thus \eqref{eq:delta} is a concrete benchmark for
a proposed arithmetic method, not a proved necessary subproblem of
every possible comparison procedure.

\section{The complete binary coefficient observer}

Work over $\Ftwo$ and set
\[
 A_d(z)=\prod_{i=1}^d(1+z^i),\qquad D=d(d+1)/2.
\]
For $b\in\{0,1\}$ let
$\Lambda_b(\sum a_nz^n)=\sum a_{2n+b}z^n$.
The state space of proper numerators is
$V_d=\{U\in\Ftwo[z]:\deg U<D\}$.

\begin{lemma}[Exact state update]\label{lem:matrix}
For $U\in V_d$,
\[
 \Lambda_b(U/A_d)=\frac{\Lambda_b(UA_d)}{A_d}.
\]
The updated numerator is again in $V_d$. In the coefficient basis its
matrix is
$T_b[i,j]=[z^{2i+b-j}]A_d$, for $0\le i,j<D$.
\end{lemma}
\begin{proof}
Since $A_d(z)^2=A_d(z^2)$ in characteristic two,
$U/A_d=(UA_d)/A_d(z^2)$. Coefficient selection commutes with a
factor in $z^2$, giving the identity. The numerator degree is at
most $\lfloor(2D-1)/2\rfloor=D-1$. Applying selection to each
monomial basis vector gives the matrix entries.
\end{proof}

\begin{theorem}[Complete selected coefficient]\label{thm:observer}
If $n=\sum_{i<L}b_i2^i$, with bits read least significant first,
then
\[
 p_d(n)\bmod2=e_0^{T}T_{b_{L-1}}\cdots T_{b_0}e_0.
\]
\end{theorem}
\begin{proof}
The initial numerator is one, represented by $e_0$. Successive
selections with $b_0,b_1,\ldots$ discard these bits of the
coefficient index. After all bits, the desired coefficient is the
constant term. As $A_d(0)=1$, this is the constant numerator
coordinate selected by $e_0^T$.
\end{proof}

This formula retains all factors and all cancellations. Boolean
reachability in its transition graph would be a different observation:
two paths contributing to one coefficient cancel in $\Ftwo$.

\begin{lemma}[All-continuation observability]\label{lem:observable}
For any field $K$, polynomial $A$ with $A(0)=1$ and $\deg A=D$,
the first $D$ coefficients of $U/A$ determine every $U$ of degree
less than $D$. For the binary observer above, distinct numerator
states are distinguished by some continuation of length at most
$\lceil\log_2 D\rceil$.
\end{lemma}
\begin{proof}
Multiplication by $1/A=1+a_1z+\cdots$ sends the coefficient
vector of $U$ to the first $D$ output coefficients by a lower
triangular matrix with diagonal one. It is invertible. Equivalently,
the first nonzero coefficient of a difference $U-V$ remains its
first nonzero coefficient after division by $A$. Its index $k<D$
has at most $\lceil\log_2D\rceil$ bits; those bits form a
distinguishing continuation. Index zero permits the empty word.
\end{proof}

This is the linear observability viewpoint familiar from weighted
automata \cite{kiefer}. It does not say that every distinguishing
continuation is available from the genuine source parameters.
A quotient preserving only the one prescribed source word remains a
different question.

\section{Differential Cartier and the surviving residues}

It is useful to normalize the same selected coefficient as a differential.
Over a perfect field $K$ of characteristic two, every rational function
has a unique expression $a(z)^2+zb(z)^2$. Define
$\mathcal C((a^2+zb^2)\,dz)=b\,dz$.
This is the differential Cartier operator. Multiplication by a square
pulls out its square root.

\begin{lemma}[Local rule]\label{lem:local}
For a fixed pole $\zeta$ and $b\in K$,
\[
 \mathcal C\!\left(\frac{b\,dz}{(z-\zeta)^{2j}}\right)=0,\qquad
 \mathcal C\!\left(\frac{b\,dz}{(z-\zeta)^{2j+1}}\right)
       =\frac{b^{1/2}\,dz}{(z-\zeta)^{j+1}}.
\]
The same definition gives
$\mathcal C(\sum a_nz^n\,dz)=\sum a_{2n+1}^{1/2}z^n\,dz$.
\end{lemma}
\begin{proof}
An even power of $h=z-\zeta$ is a square and contributes zero.
For odd powers, write $h^{2k+1}=(h^k)^2(z-\zeta)$.
The constant differential $\zeta\,dz$ has Cartier image zero,
whereas $\mathcal C(z\,dz)=dz$. Pulling out the square proves
the rule for all integer $k$, including negative $k$.
The coefficient rule follows by the same even/odd decomposition.
\end{proof}

For $e\ge1$ and $2^\ell\ge e$,
\begin{equation}\label{eq:selection}
 [z^0\,dz]\,\mathcal C^\ell\!
             \left(\frac{z^{e-1}}{A_d(z)}\,dz\right)
                   =p_d(2^\ell-e)\pmod2.
\end{equation}
Indeed $\ell$ odd selections choose original exponent $2^\ell-1$,
which in $z^{e-1}/A_d$ has coefficient $p_d(2^\ell-e)$.

\begin{lemma}[Removal of every higher pole]\label{lem:poles}
Put $\mu=\sum_{i=1}^d2^{\nu_2(i)}$ and
$a=\lceil\log_2\mu\rceil$. After $a$ Cartier applications to
a proper differential with denominator $A_d$, all surviving finite
poles are simple. Its simple residues subsequently undergo inverse
Frobenius at their fixed poles. No polynomial part is created.
\end{lemma}
\begin{proof}
The multiplicity at a root of odd order $q$ is
$\sum_{q\mid i,\ i\le d}2^{\nu_2(i)}\le\mu$.
By Lemma~\ref{lem:local}, a term of order $k>1$ survives $a$
applications only if $k-1$ is a positive multiple of $2^a$.
This is impossible for $k\le\mu\le2^a$. A simple residue is
replaced by its square root. The proper partial-fraction expansion
has no polynomial part, and every term in its Cartier image is again
proper, proving the final assertion.
\end{proof}

Here $\nu_2(i)$ denotes the exponent of two in the positive integer
$i$. The simple-pole tail is indeed eventually periodic, rather than
merely subject to a conditional period bound: all its finitely many
poles and residues lie in one finite extension of $\Ftwo$, and
inverse Frobenius is a permutation of that finite field. This
establishes a finite cycle for every supplied cutoff, without giving
a bound independent of the cutoff.

For the linked values in \eqref{eq:delta}, $C=2^\ell$ with
$\ell\ge B+1$ by \eqref{eq:growth}. Since
$\mu\le d(1+\lfloor\log_2 d\rfloor)\le B(1+\log_2B)$,
we have $a\le B$ for $B\ge16$. Thus the actual selected observation
is beyond the higher-pole transient. This pays an actual parameter
bound; it is not an unspecified promise that enough iterations occur.

The complete state can nevertheless have a long simple-pole cycle.
Pole count alone cannot prove that. For example, $dS/S$ for a
squarefree polynomial $S$ has residue one at every simple pole and
is Cartier-fixed, however many poles $S$ has.

\section{An exact period on the actual linked inputs}

Retain $u,B,C,e,d$ from \eqref{eq:delta}, and write
\[
 B=2^K,\qquad r=K-1,\qquad q=2^r+1=B/2+1,\qquad
 m_0=u-2,\qquad \Delta=m_0(m_0+1)/2.
\]
The integer $q$ need not be prime. All primitive $q$-th roots
in an algebraic closure of $\Ftwo$ will be used.

\begin{lemma}[The actual residue formula]\label{lem:residue}
The form $\omega_e=z^{e-1}A_d(z)^{-1}\,dz$ has a simple nonzero
residue at each primitive $q$-th root $\zeta$, given by
\[
 b_\zeta=\frac{\zeta^e}{A_{d-q}(\zeta)}
                 =\zeta^{e-\Delta}A_{m_0}(\zeta).
\]
\end{lemma}
\begin{proof}
The hierarchy bounds give $d/2<q\le d<2q$ and $m_0<q$.
Thus only $1+z^q$ among the factors of $A_d$ vanishes at
$\zeta$, and its derivative there is $\zeta^{-1}$.
The product $\prod_{i=1}^{q-1}(1+\zeta^i)$ is one: it is
the derivative at one of $X^q-1$, whose value is $q=1$ in
characteristic two. It follows that the residue is the first
displayed expression.

Now $d-q=q-1-m_0$. The last $m_0$ terms in that complete
product equal
$\prod_{j=1}^{m_0}(1+\zeta^{-j})
  =\zeta^{-\Delta}A_{m_0}(\zeta)$.
Taking the complementary product proves the second expression.
None of $1+\zeta^j$ for $1\le j\le m_0$ vanishes.
\end{proof}

\begin{lemma}[The bounds needed for the root count]\label{lem:bounds}
For the actual linked inputs, $r\ge96$, $\Delta<r$, and
\[
 0<\Delta-2e<q,\qquad
 \varphi(q)>\frac{2^r}{r+1}>2^{2r/3}\Delta.
\]
\end{lemma}
\begin{proof}
For $s=u+100$, $F_5(s)\ge F_1(s)\ge2^{2^s}$, so
$K\ge2^{u+100}$ and $r\ge2^{u+100}-1$.
These inequalities imply $r\ge96$ and
$\Delta=(u-2)(u-1)/2<r$.
As $e=u-4=m_0-2$,
$\Delta-2e=(m_0^2-3m_0+8)/2$ is positive and smaller
than $q$ by the same growth bound.

If an odd integer $q$ has $k$ distinct prime factors
$p_1<\cdots<p_k$, then $p_i\ge i+2$, hence
$q/\varphi(q)=\prod p_i/(p_i-1)\le k+1$.
Here $3^k\le q=2^r+1$ implies $k\le r$, so
$\varphi(q)\ge q/(r+1)>2^r/(r+1)$.
Finally $2^{r/3}>r^3$ at $r=96$ and thereafter:
the ratio of the exponential to $r^3$ increases for $r\ge96$.
Since $\Delta<r$ and $r(r+1)<r^3$, this gives
$2^r/(r+1)>2^{2r/3}\Delta$.
\end{proof}

\begin{theorem}[Exact actual residue-tuple period]\label{thm:actual-period}
The tuple $(b_\zeta)_{\zeta\ {\rm primitive}\ q}$ has exact
Frobenius period $2r$. Every eventual period of the complete
rational differential Cartier orbit of $\omega_e$ is divisible
by $2r=2(\log_2 B-1)$.
\end{theorem}
\begin{proof}
The order of two modulo $q=2^r+1$ is $2r$.
It divides $2r$ because $2^r\equiv-1\pmod q$.
Any proper divisor of $2r$ is at most $r$, and for
$s\le r$ the positive integer $2^s-1$ is smaller than $q$.
Thus no smaller order is possible. The primitive roots and their
residues belong to $\mathbb F_{2^{2r}}$, so the tuple period
divides $2r$.

Reciprocity gives
$A_{m_0}(\zeta^{-1})=\zeta^{-\Delta}A_{m_0}(\zeta)$.
If Frobenius power $r$ fixed the tuple,
Lemma~\ref{lem:residue} would give
$\zeta^{\Delta-2e}=1$ at a primitive root, contrary to
$0<\Delta-2e<q$.

Any other proper divisor $s$ of $2r$ is at most $2r/3$:
the integer quotient $(2r)/s$ is then at least three.
Fixedness at all primitive roots would make every such root a zero of
\[
 P_s(X)=A_{m_0}(X)^{2^s}
       +X^{(\Delta-e)(2^s-1)}A_{m_0}(X).
\]
This polynomial has constant coefficient one and degree exactly
$2^s\Delta$, since $e>0$ makes the second summand's degree
strictly smaller. It cannot have all $\varphi(q)$ primitive roots
as zeros, by Lemma~\ref{lem:bounds}. Thus no proper divisor is a
tuple period.

The chosen finite poles were already simple initially.
By Lemma~\ref{lem:local}, the tuple in the $j$-th Cartier state is
$(b_\zeta^{\,2^{-j}})_\zeta$ with poles fixed. If states at
$j$ and $j+p$ agree, these tuples agree. Inverse Frobenius is
bijective and its tuple orbit has period $2r$, so $2r\mid p$.
This applies also after the higher-pole transient.
\end{proof}

The proof is a tuple argument. It does not assert that every individual
residue has degree $2r$, and it does not calculate the selected
constant coefficient. In particular, the growing tuple period does
not prove that the sequence $v\mapsto\delta_4(v)$ is nonperiodic.

\section{What has and has not been reduced}

For each actual input, the observer is finite and complete. Its size
grows with $d$, and its exact rational tail has no uniformly bounded
state period. These facts explain why a proposed solution must do more
than attach a finite automaton to each supplied cutoff.

One possible improvement would preserve only the scalar observations
allowed by the linked source inputs. The all-continuation
Lemma~\ref{lem:observable} leaves that possibility open: it distinguishes
states using arbitrary binary continuations, not just the prescribed
source word. Another possibility would classify complete descent-count
residues without first evaluating the isolated boundary bit. Neither
possibility is supplied by the present arguments.

There are therefore three distinct mathematical outcomes here:
an exact congruence transferring a growing-budget count to a
restricted partition; a complete finite observer for that coefficient;
and an exact full-state period obstruction on genuine inputs.
Their proofs do not settle the remaining fixed-control all-ground
canonical comparison problem. The source-expressivity companion
\cite{wwexpressivity} explains why complete-count attainment is an
actual necessary arithmetic problem, while extracting one summand
from the complete count requires an additional interpretation.

\begin{thebibliography}{99}
\bibitem{ww}Gunnar Wilken and Andreas Weiermann.
Derivation lengths classification of G\"odel's T extending Howard's assignment.
\emph{Logical Methods in Computer Science} \textbf{8}(1:19) (2012), 1--44.
\href{https://doi.org/10.2168/LMCS-8(1:19)2012}
{doi:10.2168/LMCS-8(1:19)2012}.
\href{https://lmcs.episciences.org/1073}{Official article and full text}.
The 2012 revised article is the version used throughout.

\bibitem{bcw}W.~Buchholz, A.~Cichon and A.~Weiermann.
\emph{A Uniform Approach to Fundamental Sequences and Hierarchies}.
Mathematical Logic Quarterly \textbf{40}(2) (1994), 273--286.
\href{https://doi.org/10.1002/malq.19940400212}{doi:10.1002/malq.19940400212}.
The norm-defined predecessor framework is discussed on p.~274; the
additive symbol-count norm appears on p.~283.

\bibitem{kiefer}Stefan Kiefer.
\emph{Notes on equivalence and minimization of weighted automata}.
arXiv:2009.01217v1, 2 September 2020.
\href{https://doi.org/10.48550/arXiv.2009.01217}
{doi:10.48550/arXiv.2009.01217}.
The supplied first version is the version inspected.

\bibitem{wwfoundations}\emph{The Wilken--Weiermann Comparison Problem: Foundations and a Guide to the Results}. Companion working manuscript, 2026.\
\href{../05-wilken-weiermann-foundations/paper.pdf}{Included manuscript}.

\bibitem{wwexpressivity}\emph{Actual-Source Arithmetic and Undecidability Boundaries for Canonical Comparison}. Companion working manuscript, 2026.\
\href{../09-canonical-source-expressivity/paper.pdf}{Included manuscript}.
\end{thebibliography}

\end{document}
